\documentclass[a4paper, serif, 10pt]{moderncv}

%% moderncv
% color and style
\moderncvcolor{black}
\moderncvstyle{banking}

%% page layout/margins
\usepackage[top=2.5cm, bottom=2.5cm, left=1.5cm, right=1.5cm]{geometry}

\nopagenumbers{}

%% Babel config for Dutch language
% ref:
% - https://latex3.github.io/babel/guides/locale-dutch.html
\usepackage[dutch]{babel}

%% fontspec
\usepackage{fontspec}

\IfFontExistsTF{Linux Libertine}{
  \setmainfont[Ligatures={TeX, Common}, Numbers=OldStyle]{Linux Libertine}
}{}
\IfFontExistsTF{Linux Libertine O}{
  \setmainfont[Ligatures={TeX, Common}, Numbers=OldStyle]{Linux Libertine O}
}{}

%% CTeX
% CJK support, used to show CJK characters in the resume
%
% - heading=false: disable ctex heading
% - punct=kaiming: use kaiming punctuations styles for CJK
% - scheme=plain: use plain scheme, do not override \normalsize font size
% - space=auto: space settings for CJK characters
%
% CJK fonts are resolved per locale language, e.g. Noto Serif/Sans CJK JP
% for Japanese resumes, with fallback to the next region if not installed
%
% ref:
% - http://ctan.mirrorcatalogs.com/language/chinese/ctex/ctex.pdf
\usepackage[UTF8, heading=false, punct=kaiming, scheme=plain, space=auto]{ctex}

\IfFontExistsTF{Noto Serif CJK SC}{
  \setCJKmainfont{Noto Serif CJK SC}
}{}
\IfFontExistsTF{Noto Sans CJK SC}{
  \setCJKsansfont{Noto Sans CJK SC}
}{}

%% line spacing
\usepackage{setspace}
\setstretch{1.125}

%% URL styling - use normal text instead of monospace
\urlstyle{same}

%% Patch \httplink and \httpslink to handle full URLs
% The original moderncv macros always prepend http:// or https:// to URLs.
% This causes issues when the URL already has a protocol (e.g., https://example.com)
% resulting in malformed URLs like https://https://example.com.
% This patch checks if the URL already contains "://" and uses it directly if so.
% Note: We use \str_set:Nx (x-expansion) to expand macros like \@homepage first.
\ExplSyntaxOn
\str_new:N \g_yamlresume_url_str

\RenewDocumentCommand{\httpslink}{O{}m}{
  \str_set:Nx \g_yamlresume_url_str {#2}
  \str_if_in:NnTF \g_yamlresume_url_str {://}
    { \url{#2} }
    { \url{https://#2} }
}

\RenewDocumentCommand{\httplink}{O{}m}{
  \str_set:Nx \g_yamlresume_url_str {#2}
  \str_if_in:NnTF \g_yamlresume_url_str {://}
    { \url{#2} }
    { \url{http://#2} }
}
\ExplSyntaxOff

\name{Sanne de Vries}{}
\title{AI-engineer | Machine Learning \& MLOps}
\phone[mobile]{+31 6 1234 5678}
\email{sanne@sannedevries.nl}
\homepage{https://www.sannedevries.nl}

\address{Keizersgracht 123, Amsterdam, Noord-Holland, Nederland, 1015 CJ}{}{}

\extrainfo{{\small \faLinkedin}\ \href{https://www.linkedin.com/in/sannedevries-ai}{@sannedevries-ai} {} {} {} • {} {} {} 
{\small \faGithub}\ \href{https://github.com/sannedevries}{@sannedevries} {} {} {} • {} {} {} 
{\small \faStackOverflow}\ \href{https://stackoverflow.com/users/18422031/sanne-de-vries}{@sanne-de-vries}}

\begin{document}

\maketitle

\section{Basis Informatie}

\cvline{}{AI-engineer met 6+ jaar ervaring in het ontwerpen, bouwen en opschalen van machine learning-oplossingen van prototype tot productie. Sterk in het vertalen van businessvragen naar meetbare modellen en het neerzetten van betrouwbare MLOps-pipelines.
%
\begin{itemize}
\item Ervaren in Python, PyTorch, TensorFlow en cloudplatformen (AWS, Azure, GCP)
\item Combineert kennis van natuurlijke taalverwerking en computer vision met pragmatische software-engineering
\item Gedreven teamspeler die graag kennis deelt en multidisciplinaire teams meeneemt in verantwoorde AI
\end{itemize}}
\section{Opleidingen}

\cventry{sep 2017–jul 2019}
        {Master, Computerwetenschappen - Kunstmatige Intelligentie, Score: 8.2/10}
        {Technische Universiteit Delft}
        {\url{https://www.tudelft.nl}}
        {}
        {\begin{itemize}
\item Gespecialiseerd in machine learning, natuurlijke taalverwerking en uitlegbaarheid van modellen
\item Afstudeerproject over transparante deep learning-modellen voor klinische besluitvorming, beoordeeld met een 9,0
\item Praktijkervaring opgedaan via onderzoekssamenwerkingen met twee industriële partners
\item Vakken op het snijvlak van algoritmiek, statistiek en verantwoorde AI
\end{itemize}
\textbf{Cursussen}: Machine Learning, Deep Learning, Natuurlijke taalverwerking, Computer Vision, Reinforcement Learning, MLOps en Cloud Computing, Ethiek van AI, Data Engineering}

\cventry{sep 2014–jun 2017}
        {Bachelor, Informatiekunde, Score: 7.8/10}
        {Universiteit Utrecht}
        {\url{https://www.uu.nl}}
        {}
        {\begin{itemize}
\item Brede basis in informatiesystemen, data-analyse en softwareontwikkeling
\item Minor gevolgd in toegepaste statistiek en datawetenschap
\item Scriptie over predictive analytics voor openbare voorzieningen, afgerond met een 8,5
\end{itemize}
\textbf{Cursussen}: Programmeren in Python, Databases en SQL, Statistiek en Kansrekening, Algoritmen en Datastructuren, Informatiesystemen, Inleiding Machine Learning}
\section{Werkervaring}

\cventry{mrt 2022–Heden}
        {Senior AI-engineer}
        {Aiden Analytics}
        {\url{https://www.aidenanalytics.nl}}
        {}
        {\begin{itemize}
\item Leidt een team van vier engineers dat generatieve AI-toepassingen bouwt voor klanten in de financiële sector en de zorg
\item Ontwierp een retrieval-augmented generatiearchitectuur die de responstijd met 45\% verlaagde en de feitelijke nauwkeurigheid meetbaar verbeterde
\item Zette een MLOps-platform op met CI/CD, modelregistratie, monitoring en drift-detectie, waardoor releases van weken naar dagen gingen
\item Stelde interne richtlijnen op voor verantwoorde AI, inclusief bias-testen, modelkaarten en menselijke controle bij hoog-risico beslissingen
\item Coacht junior engineers en verzorgt kennissessies over LLM-evaluatie en prompt-engineering
\end{itemize}
\textbf{Trefwoorden}: MLOps, Generatieve AI, LLM-evaluatie, Kubernetes, Verantwoorde AI}

\cventry{sep 2019–feb 2022}
        {Machine Learning Engineer}
        {MediQ Data}
        {\url{https://www.mediqdata.nl}}
        {}
        {\begin{itemize}
\item Ontwikkelde en productiseerde voorspelmodellen voor patiëntstromen, met een reductie van 18\% in gemiste afspraken
\item Bouwde NLP-pipelines voor het anonimiseren en structureren van vrije tekst uit elektronische patiëntendossiers
\item Automatiseerde het trainingsproces met PyTorch en Airflow, inclusief reproduceerbare experimenten en versiebeheer van datasets
\item Werkte samen met data-ethici en juristen om te voldoen aan de AVG en medische richtlijnen
\item Presenteerde resultaten aan niet-technische stakeholders en vertaalde die naar concrete productkeuzes
\end{itemize}
\textbf{Trefwoorden}: PyTorch, NLP, Airflow, AVG, Voorspelmodellen}

\cventry{aug 2018–aug 2019}
        {Data Scientist}
        {Blauwdruk Digital}
        {\url{https://www.blauwdrukdigital.nl}}
        {}
        {\begin{itemize}
\item Analyseerde klantgedrag voor e-commerceklanten en bouwde aanbevelingsmodellen die de conversie met 12\% verhoogden
\item Ontwikkelde herbruikbare feature-engineering-bibliotheek in Python die door drie teams werd gebruikt
\item Realiseerde dashboards en rapportages in samenwerking met marketing- en productteams
\item Verkende en vergeleek modellen met A/B-testen en statistische validatie
\end{itemize}
\textbf{Trefwoorden}: Python, Aanbevelingssystemen, A/B-testen, SQL, Feature Engineering}
\section{Talen}

\cvline{Nederlands}{Moedertaal of tweetalige vaardigheid \hfill \textbf{Trefwoorden}: Moedertaal}
\cvline{Engels}{Volledige professionele vaardigheid \hfill \textbf{Trefwoorden}: Dagelijks gebruikt op het werk, Technische documentatie}
\cvline{Duits}{Beperkte werkvaardigheid \hfill \textbf{Trefwoorden}: Basale gesprekken}
\section{Vaardigheden}

\cvline{Machine Learning}{Expert \hfill \textbf{Trefwoorden}: PyTorch, TensorFlow, Scikit-learn, XGBoost, Hugging Face}
\cvline{MLOps en Infrastructuur}{Geavanceerd \hfill \textbf{Trefwoorden}: Docker, Kubernetes, MLflow, Airflow, Terraform}
\cvline{Data Engineering}{Geavanceerd \hfill \textbf{Trefwoorden}: Python, SQL, Spark, dbt, Kafka}
\cvline{Cloud}{Geavanceerd \hfill \textbf{Trefwoorden}: AWS SageMaker, Azure ML, Google Vertex AI}
\cvline{AI-ethiek en Governance}{Intermediair \hfill \textbf{Trefwoorden}: Modelkaarten, Bias-detectie, Uitlegbaarheid, AVG}
\section{Onderscheidingen}

\cventry{nov 2023}
        {Dutch Data Science Awards}
        {Beste AI-project van het jaar}
        {}
        {}
        {Toegekend voor een retrieval-augmented generatieoplossing die klantmedewerkers ondersteunt bij complexe vragen, met aantoonbaar hogere nauwkeurigheid en kortere afhandeltijd.}

\cventry{jul 2019}
        {Technische Universiteit Delft}
        {Cum laude afstudeerproject}
        {}
        {}
        {Scriptie over uitlegbaarheid van deep learning in de zorg beoordeeld met een 9,0 en genomineerd voor de facultaire scriptieprijs.}
\section{Certificaten}

\cventry{apr 2023}
        {Amazon Web Services}
        {AWS Certified Machine Learning - Specialty}
        {\url{https://aws.amazon.com/certification/}}
        {}
        {}

\cventry{jun 2022}
        {Google Cloud}
        {Professional Machine Learning Engineer}
        {\url{https://cloud.google.com/learn/certification}}
        {}
        {}

\cventry{feb 2021}
        {DeepLearning.AI}
        {Deep Learning Specialization}
        {\url{https://www.coursera.org/specializations/deep-learning}}
        {}
        {}
\section{Publicaties}

\cventry{mrt 2020}
        {Uitlegbare voorspelmodellen voor klinische besluitvorming: een praktijkstudie}
        {Journal of Biomedical Informatics}
        {\url{https://www.sciencedirect.com/journal/journal-of-biomedical-informatics}}
        {}
        {\begin{itemize}
\item Onderzoekt hoe attention-visualisaties en counterfactual explanations clinici ondersteunen bij het interpreteren van modeluitkomsten
\item Beschrijft een evaluatie met 42 clinici in twee ziekenhuizen
\item Geeft praktische richtlijnen voor het documenteren en valideren van klinische modellen
\end{itemize}}
\section{Referenties}

\cventry{\emaillink[m.vandenberg@tudelft.nl]{m.vandenberg@tudelft.nl}}
        {Hoogleraar Kunstmatige Intelligentie}
        {Dr. Marijke van den Berg}
        {+31 15 278 1234}
        {}
        {Sanne combineert sterke technische vaardigheden met een scherp oog voor de maatschappelijke impact van AI. Ze is zelfstandig, communiceert helder met niet-technische stakeholders en tilt teams naar een hoger niveau.}
\section{Projecten}

\cventry{jan 2023–Heden}
        {Open-source assistent die interne documentatie doorzoekbaar maakt met retrieval-augmented generatie}
        {Vraagbaak}
        {\url{https://github.com/sannedevries/vraagbaak}}
        {}
        {\begin{itemize}
\item Combineert een vector-database met een taalmodel voor betrouwbare antwoorden met bronverwijzingen
\item Inclusief evaluatieharness voor hallucinatiedetectie en antwoordkwaliteit
\item Wordt door drie organisaties in productie gebruikt
\end{itemize}
\textbf{Trefwoorden}: RAG, Vector Database, LLM, Open Source}

\cventry{mei 2021–dec 2022}
        {Computer vision-pipeline voor het automatisch classificeren van inspectiefoto's in de maakindustrie}
        {Beeldscan}
        {\url{https://github.com/sannedevries/beeldscan}}
        {}
        {\begin{itemize}
\item Transfer learning met een CNN backbone, getraind op een gelabelde dataset van 120.000 afbeeldingen
\item Bereikte een gemiddelde precisie van 96\% bij vier schadecategorieën
\item Geïmplementeerd als containerized inference-service met monitoring op datadrift
\end{itemize}
\textbf{Trefwoorden}: Computer Vision, Transfer Learning, Docker}
\section{Interesses}

\cvline{Sport}{Hardlopen, Wielrennen, Zwemmen}
\cvline{Muziek}{Piano, Jazz}
\cvline{Overig}{Schaken, Open Source, Filosofie van AI}
\section{Vrijwilligerswerk}

\cventry{sep 2020–Heden}
        {Mentor datawetenschap}
        {Stichting AI voor Iedereen}
        {\url{https://www.aivooriedereen.nl}}
        {}
        {\begin{itemize}
\item Begeleidt jaarlijks vier carrièreswitchers bij hun eerste stappen in machine learning
\item Verzorgt gastlessen over AI-geletterdheid op middelbare scholen
\item Organiseert maandelijkse bijeenkomsten voor de open-source community
\end{itemize}}
    

\cventry{feb 2022–Heden}
        {Vrijwilliger werkgroep Ethiek}
        {NL AI Coalitie}
        {\url{https://www.nlaic.nl}}
        {}
        {\begin{itemize}
\item Draagt bij aan publieke richtlijnen voor transparante en uitlegbare AI-toepassingen
\item Reviewt praktijkcases van organisaties die hun AI-systemen willen documenteren
\end{itemize}}
    
\end{document}